\documentclass[11pt,letterpaper]{article}
\usepackage[margin=0.75in]{geometry}
\usepackage{mathptmx}
\usepackage{booktabs}
\usepackage{amssymb}
\setlength{\parindent}{0pt}
\setlength{\parskip}{5pt}
\pagestyle{empty}

\begin{document}

\begin{center}
    {\large \textbf{CERTIFICATION OF MICRO ENTITY STATUS -- GROSS INCOME BASIS}}\\[2pt]
    {\small \textbf{Pursuant to 35 U.S.C.\ \S\ 123(a) and 37 CFR 1.29 (Equivalent to USPTO Form PTO/SB/15A)}}
\end{center}
\vspace{4pt}

\begin{tabular}{|p{2.2in}|p{4.5in}|}
\hline
\textbf{Application Number / Filing Date} & Provisional Application filed herewith (November 5, 2026) \\
\hline
\textbf{First Named Inventor} & Dionisio Chiuratto Agourakis \\
\hline
\textbf{Title of Invention} & APPARATUS, SYSTEMS, AND METHODS FOR NATIVE 16-DIMENSIONAL SEDENION SIGNAL PROCESSING, PRE-ADC ANALOG ZERO-DIVISOR INTERFERENCE CANCELLATION, SINGLE-SNAPSHOT NON-ASSOCIATIVE CURVATURE SPOOFING DETECTION, AND ZERO-DIVISOR KEYED MODULATION \\
\hline
\end{tabular}

\vspace{8pt}
\textbf{The applicant hereby certifies the following under 35 U.S.C.\ \S\ 123(a) and 37 CFR 1.29(a):}

\begin{enumerate}
    \item \textbf{Small Entity Qualification (37 CFR 1.29(a)(1)):} The applicant qualifies as a small entity as defined in 37 CFR 1.27 without relying on a government use license exception under 37 CFR 1.27(a)(4).
    \item \textbf{Application Filing Limit (37 CFR 1.29(a)(2)):} Neither the applicant nor the inventor nor a joint inventor has been named as the inventor or a joint inventor on more than four previously filed U.S.\ patent applications, excluding provisional applications, international applications for which the basic national fee under 35 U.S.C.\ \S\ 41(a) was not paid, and applications assigned to a previous employer as a result of previous employment.
    \item \textbf{Gross Income Limit (37 CFR 1.29(a)(3)):} Neither the applicant nor the inventor nor a joint inventor, in the calendar year preceding the calendar year in which the applicable fee is being paid, had a gross income, as defined in section 61(a) of the Internal Revenue Code of 1986 (or equivalent foreign gross income converted to U.S.\ dollars at the applicable exchange rate under 37 CFR 1.29(c)), exceeding the ``Maximum Qualifying Gross Income'' reported on the USPTO website (three times the median household income for that preceding calendar year, as most recently reported by the Bureau of the Census).
    \item \textbf{No Assignment/Obligation to Non-Qualifying Entity (37 CFR 1.29(a)(4)):} Neither the applicant nor the inventor nor a joint inventor has assigned, granted, or conveyed, nor is under an obligation by contract or law to assign, grant, or convey, a license or other ownership interest in the application concerned to an entity that, in the calendar year preceding the calendar year in which the applicable fee is being paid, had a gross income exceeding the ``Maximum Qualifying Gross Income'' limit described in item 3 above.
\end{enumerate}

\vspace{8pt}
\textbf{SIGNATURE OF APPLICANT / SOLE INVENTOR (37 CFR 1.33(b) \& 37 CFR 1.4(d)(2)):}\\[4pt]
\begin{tabular}{|p{2.6in}|p{2.3in}|p{1.6in}|}
\hline
\textbf{Signature (37 CFR 1.4(d)(2))} & \textbf{Typed or Printed Name} & \textbf{Date} \\
\hline
\rule{0pt}{18pt}\texttt{/Dionisio Chiuratto Agourakis/} & Dionisio Chiuratto Agourakis & November 5, 2026 \\
\hline
\textbf{Telephone:} +55 11 98787-4110 & \textbf{Email:} \texttt{dionisio.ca@gmail.com} & \textbf{Capacity:} Sole Inventor \\
\hline
\end{tabular}

\end{document}
